\documentclass[11pt,a4paper]{article}
\usepackage[T1]{fontenc}
\usepackage[utf8]{inputenc}
\usepackage[spanish]{babel}
\usepackage{amsmath,booktabs,geometry,hyperref}
\geometry{margin=24mm}
\title{Informe piloto: bigramas sobre bitácoras sintéticas}
\author{Laboratorio S01 · ejemplo ejecutado en Python}
\date{17 de octubre de 2026}
\begin{document}\maketitle
\begin{abstract}
Se demuestra el protocolo mínimo para comparar modelos probabilísticos con
métricas homogéneas y un informe reproducible. Los documentos se han
fabricado para la clase: la experiencia no constituye evidencia de campo.
\end{abstract}
\section{Pregunta y datos}
Se ajustaron dos bigramas de caracteres sobre 8 textos generales y 8 textos
técnicos, respectivamente. Se evaluaron en 6 textos técnicos distintos,
igualmente sintéticos. No se descargaron corpus ni se ejecutaron modelos LLM.
\section{Método}
Vocabulario común fijado con entrenamiento, minúsculas, recorte de espacios
extremos, suavizado aditivo $\alpha=1$, BOS y EOS declarados. La métrica es
$\mathrm{BPC}(z)=-\log_2\widehat P(z,\mathrm{EOS})/|z|$,
contando puntos de código Unicode sin EOS en el denominador.
\section{Resultados reproducidos}
\begin{center}\begin{tabular}{lrrrr}\toprule
Documento & Caracteres & BPC general & BPC técnico & Diferencia\\\midrule
E01 & 27 & 4.167 & 3.609 & -0.559 \\
E02 & 34 & 4.269 & 3.448 & -0.821 \\
E03 & 32 & 3.900 & 3.533 & -0.367 \\
E04 & 37 & 4.188 & 3.665 & -0.523 \\
E05 & 36 & 4.176 & 3.558 & -0.619 \\
E06 & 31 & 3.840 & 3.487 & -0.353 \\
\bottomrule\end{tabular}\end{center}

\noindent Diferencia media (técnico menos general): $-0.54024842$ bits/carácter. El intervalo percentil al 95\% fue
$[-0.66770598,-0.42038147]$, obtenido por enumerar
$6^6=46656$ remuestras de documentos con reemplazo y conservar el pareamiento.
\section{Microexperimento BPE de tokenizacion (IC10)}
Ajustamos dos tokenizadores BPE didacticos, uno en texto general y otro
tecnico, con 20 fusiones cada uno; las reglas se aprendieron \emph{solo}
de entrenamiento y se aplicaron a los mismos seis documentos reservados:
\begin{center}\begin{tabular}{lrr}\toprule
Tokenizador & Tokens & Tokens/palabra\\\midrule
BPE general & 148 & 4.228571\\
BPE tecnico & 127 & 3.628571\\\bottomrule
\end{tabular}\end{center}
Son 35 palabras, diferencia total $-21$ tokens. Es un resultado de
\textbf{longitud tokenizada}, no de prediccion ni de un LLM.
\section{Interpretación y límites}
La comparación favorece al modelo técnico para \emph{estos seis ejemplos
inventados}. No estima efectos poblacionales; el intervalo tampoco controla
sesgos de selección o de redacción. No se han comparado tokenizadores modernos.
En un estudio real se exigen licencia de corpus, particiones blindadas,
baselines, varias semillas y validación externa.
\section{Vínculo con investigación y objetivo Q1}
El problema IC10 plantea estudiar la penalización de tokenización en textos
técnicos en español. La experiencia de esta sesión solo enseña a medir y
documentar. Un artículo candidato a una revista Q1 exige contribución
novedosa, literatura crítica, evidencia sobre datos reales, robustez y
reproducibilidad. La obligación de S01 es identificar dos áreas de interés;
la selección formal del problema ocurre en S05.
\end{document}
